\chapter{Lowest Common Ancestor and Range Structures}

% --- chapter-local shorthands (do not clash with common.tex) ---
\newcommand{\dep}{\operatorname{dep}}
\newcommand{\pa}{\operatorname{par}}
\newcommand{\up}{\mathrm{up}}
\newcommand{\spt}{\mathrm{sp}}
\newcommand{\eu}{\mathrm{E}}
\newcommand{\de}{\mathrm{H}}
\newcommand{\fst}{\mathrm{first}}
\newcommand{\lz}{\mathrm{lazy}}
\providecommand{\lowbit}{\operatorname{lowbit}}

Hang a tree from its root and let it dangle. Pick two nodes $u$ and $v$ down in the
foliage. Trace the path from each up toward the root. The two paths run separately for a
while, then meet, and from the meeting point upward they coincide. That first shared node --
the deepest node that is an ancestor of both -- is the \emph{lowest common ancestor} of $u$
and $v$, written $\operatorname{lca}(u,v)$. It is the joint where their family histories
merge.

The LCA is the single most useful primitive on rooted trees. The distance between two nodes
in a weighted tree is $\dep(u)+\dep(v)-2\dep(\ell)$ where $\ell=\operatorname{lca}(u,v)$.
``Is $u$ an ancestor of $v$?'' is just ``$\operatorname{lca}(u,v)=u$?''. Path queries,
auto-completion tries, suffix-tree string matching, phylogenetics, and half the tree
problems in competitive programming bottom out in an LCA call. So we want it fast: after
linear or near-linear preprocessing, answer each query in $\Oh(\log n)$ or even $\Oh(1)$
time.

This chapter tells two intertwined stories. The first is LCA itself, by two routes:
\emph{binary lifting}, which precomputes ancestor pointers at power-of-two heights and
climbs in $\Oh(\log n)$ per query; and the \emph{Euler tour} reduction, which flattens the
tree into an array and turns each LCA query into a \emph{range-minimum query} (RMQ) over
that array. The second story is the toolbox of \emph{range structures} -- sparse tables,
segment trees, and Fenwick (binary indexed) trees -- that answer ``summarize this contiguous
block'' queries. RMQ is one such question; sums, maxima, and assignment-on-ranges are
others. These structures are the workhorses underneath nearly every tree algorithm in this
part of the book, so we build them carefully and prove their bounds. The chapter closes with
the elegant fact that LCA and RMQ are \emph{the same problem} up to linear-time
transformations in either direction.

We trace a concrete eleven-node tree by hand, all the way through, watching every table
fill in.

\chapterline

\section{Rooted Trees, Depth, and Ancestry}
\label{ch23:setup}

Fix a tree $T=(V,E)$ with $n=\abs{V}$ vertices, rooted at a chosen vertex $r$. Rooting
orients every edge: each non-root vertex $v$ has a unique \emph{parent} $\pa(v)$, the
neighbor of $v$ on the path to $r$, and the other neighbors are its \emph{children}. The
\emph{depth} $\dep(v)$ is the number of edges on the path from $r$ to $v$, so $\dep(r)=0$.

\begin{definition}[Ancestor, descendant, LCA]
\label{ch23:def-lca}
Vertex $a$ is an \emph{ancestor} of $v$ if $a$ lies on the (unique) path from $r$ to $v$;
equivalently $a=v$ or $a$ is an ancestor of $\pa(v)$. Then $v$ is a \emph{descendant} of
$a$. The \emph{lowest common ancestor} $\operatorname{lca}(u,v)$ is the ancestor of both
$u$ and $v$ of maximum depth.
\end{definition}

\begin{intuitionbox}
Why is the LCA well defined and unique? The set of common ancestors of $u$ and $v$ always
contains the root, so it is nonempty. The ancestors of any vertex form a chain ordered by
depth (the root-to-vertex path), so the common ancestors -- being a subset of both chains --
also form a chain. A finite nonempty chain has a unique deepest element. That element is the
LCA.
\end{intuitionbox}

Throughout we run our methods on the same example so the data structures can be compared
side by side.

\begin{figure}[H]
\centering
\begin{tikzpicture}[
  every node/.style={circle,draw=nodeblue,thick,fill=blue!4,minimum size=7mm,
    inner sep=0pt,font=\small},
  level distance=12mm,sibling distance=15mm,>=Stealth]
\node (1) {1}
  child {node (2) {2}
    child {node (4) {4}}
    child {node (5) {5}
      child {node (8) {8}}
      child {node (9) {9}}
    }
  }
  child {node (3) {3}
    child {node (6) {6}
      child {node (10) {10}}
      child {node (11) {11}}
    }
    child {node (7) {7}}
  };
\node[draw=none,fill=none,right=18mm of 1,font=\footnotesize\itshape,text width=46mm]
  {root $r=1$; depths\\
   $\dep(1){=}0$;\ $\dep(2){=}\dep(3){=}1$;\\
   $\dep(4){=}\dep(5){=}\dep(6){=}\dep(7){=}2$;\\
   $\dep(8){=}\dep(9){=}\dep(10){=}\dep(11){=}3$.};
\end{tikzpicture}
\caption{The running tree $T$. For example $\operatorname{lca}(8,9)=5$,
$\operatorname{lca}(8,4)=2$, $\operatorname{lca}(10,7)=3$, and
$\operatorname{lca}(8,11)=1$.}
\label{ch23:fig-tree}
\end{figure}

We write $\pa$ for the parent array. For $T$,
\[
\pa = [\,-,\,1,\,1,\,2,\,2,\,3,\,3,\,5,\,5,\,6,\,6\,],
\]
indexed $1\ldots 11$ with $\pa(1)$ undefined (the root). Computing $\pa$ and $\dep$ is one
depth-first or breadth-first sweep, $\Oh(n)$ time. Everything below assumes these two arrays
are in hand.

\section{The Naive Climb, and Why We Want Better}
\label{ch23:naive}

The definition hands us an algorithm. To find $\operatorname{lca}(u,v)$, first lift the
deeper of the two up to the depth of the shallower by repeatedly following parent pointers;
now $u$ and $v$ sit at equal depth. Then move both up one step at a time until they collide.

\begin{algorithm}[H]
\DontPrintSemicolon
\caption{\textsc{NaiveLCA}$(u,v)$ -- climb one step at a time}
\While{$\dep(u) > \dep(v)$}{$u \gets \pa(u)$\;}
\While{$\dep(v) > \dep(u)$}{$v \gets \pa(v)$\;}
\While{$u \neq v$}{$u \gets \pa(u)$\; $v \gets \pa(v)$\;}
\Return $u$\;
\end{algorithm}

This is correct -- once the depths agree, the two climbers stay at equal depth and their
ancestor chains must meet exactly at the LCA -- but a single query costs $\Oh(n)$ on a path
graph (a ``bamboo''). With $q$ queries that is $\Oh(nq)$, hopeless for the large $n$ and $q$
these structures are meant for. The two big ideas of the chapter both attack the slow
$\Oh(n)$ climb: binary lifting climbs in jumps of $1,2,4,8,\dots$; the Euler tour avoids
climbing entirely.

\section{Binary Lifting}
\label{ch23:binary-lifting}

\begin{intuitionbox}
Climbing one parent at a time wastes the structure of the integers. Any number of steps
$k$ has a binary expansion: $13 = 8+4+1$. If we could jump up by exactly $8$, then $4$, then
$1$ ancestors in one move each, we would reach the $13$th ancestor in three moves instead of
thirteen. So precompute, for every vertex, its $2^j$-th ancestor for each
$j=0,1,\dots,\lfloor\log_2 n\rfloor$. Then \emph{any} climb of $k$ steps decomposes into at
most $\lceil\log_2 n\rceil$ power-of-two jumps.
\end{intuitionbox}

Let $\up[v][j]$ denote the $2^j$-th ancestor of $v$ (or a sentinel $0$, meaning ``above the
root'', when no such ancestor exists). The base case is $\up[v][0]=\pa(v)$. The recurrence is
the heart of the method: the $2^j$-th ancestor is the $2^{j-1}$-th ancestor \emph{of} the
$2^{j-1}$-th ancestor,
\begin{equation}
\up[v][j] \;=\; \up[\,\up[v][j-1]\,][\,j-1\,].
\label{ch23:lift-recur}
\end{equation}
Two half-jumps of length $2^{j-1}$ compose into one full jump of length $2^j$. We fill the
table column by column (all $j{=}0$ first, then all $j{=}1$, $\dots$) so that the right-hand
side of \eqref{ch23:lift-recur} is always already computed.

\begin{algorithm}[H]
\DontPrintSemicolon
\caption{\textsc{BuildLifting}$(\pa,\dep,n)$}
$L \gets \lfloor \log_2 n\rfloor$\;
\For{$v \gets 1$ \KwTo $n$}{$\up[v][0] \gets \pa(v)$ \tcp*{$\pa(r)=0$, the sentinel}}
\For{$j \gets 1$ \KwTo $L$}{
  \For{$v \gets 1$ \KwTo $n$}{
    $\up[v][j] \gets \up[\,\up[v][j-1]\,][\,j-1\,]$\;
  }
}
\Return $\up$\;
\end{algorithm}

The table has $n(L{+}1)=\Th(n\log n)$ entries, each filled in $\Oh(1)$, so preprocessing is
$\Th(n\log n)$ time and space. We adopt the convention $\up[0][j]=0$: jumping above the root
stays at the sentinel, which makes the recurrence total without special cases.

\subsection{Answering a query}

To compute $\operatorname{lca}(u,v)$: first lift the deeper vertex to the depth of the
shallower, using the binary expansion of the depth difference. Now both are at the same
depth. If they are already equal, that is the answer. Otherwise climb both \emph{in
lockstep} by the largest jumps that keep them apart, from $j=L$ down to $0$: if
$\up[u][j]\neq\up[v][j]$ the $2^j$-ancestors are still below the LCA, so jump both there;
if they are equal we would overshoot, so skip that bit. After the loop, $u$ and $v$ sit
exactly one step below the LCA, and $\pa(u)$ is the answer.

\begin{algorithm}[H]
\DontPrintSemicolon
\caption{\textsc{LiftLCA}$(u,v)$}
\If{$\dep(u) < \dep(v)$}{swap $u \leftrightarrow v$ \tcp*{make $u$ the deeper one}}
$\Delta \gets \dep(u)-\dep(v)$\;
\For{$j \gets L$ \KwTo $0$}{
  \lIf{bit $j$ of $\Delta$ is set}{$u \gets \up[u][j]$}
}
\lIf{$u = v$}{\Return $u$}
\For{$j \gets L$ \KwTo $0$}{
  \If{$\up[u][j] \neq \up[v][j]$}{
    $u \gets \up[u][j]$\; $v \gets \up[v][j]$\;
  }
}
\Return $\up[u][0]$\;
\end{algorithm}

\begin{whybox}
\textbf{Why the lockstep descent lands exactly below the LCA.} Let
$\ell=\operatorname{lca}(u,v)$ after the depths are equalized; say both vertices are at depth
$d$ and $\ell$ is at depth $d'\le d$, so $k=d-d'$ steps separate each of $u,v$ from $\ell$.
We claim the second loop drives both to depth $d'{+}1$, i.e. to the two distinct children of
$\ell$ on the respective paths.

The invariant maintained is: \emph{$u$ and $v$ stay at equal depth, and neither has reached
$\ell$ (so $u\neq v$, and $\up[u][\cdot]$, $\up[v][\cdot]$ are genuine ancestors).} A jump at
height $j$ is taken exactly when $\up[u][j]\neq\up[v][j]$. If the two $2^j$-ancestors differ
they are both strictly below $\ell$ (equal-depth proper ancestors that are not yet the LCA
cannot have met), so jumping preserves the invariant. If they are equal, that common
ancestor is at or above $\ell$, and jumping would overshoot, so we correctly skip. Processing
$j$ from high to low realizes the binary expansion of $k-1$: the loop reduces the remaining
gap to exactly $1$. Hence $u,v$ end as the two children of $\ell$, and $\up[u][0]=\ell$.
\end{whybox}

Each loop runs $L{+}1$ times with $\Oh(1)$ work, so a query is $\Oh(\log n)$.

\begin{labbox}
\textbf{Build the lifting table for $T$ (Figure~\ref{ch23:fig-tree}).} Here $n=11$ so
$L=\lfloor\log_2 11\rfloor = 3$. Column $j{=}0$ is just $\pa$. Column $j{=}1$ applies the
recurrence $\up[v][1]=\pa(\pa(v))$, and so on:
\[
\renewcommand{\arraystretch}{1.15}
\begin{array}{c|ccccccccccc}
v & 1 & 2 & 3 & 4 & 5 & 6 & 7 & 8 & 9 & 10 & 11\\\hline
\up[v][0] & 0 & 1 & 1 & 2 & 2 & 3 & 3 & 5 & 5 & 6 & 6\\
\up[v][1] & 0 & 0 & 0 & 1 & 1 & 1 & 1 & 2 & 2 & 3 & 3\\
\up[v][2] & 0 & 0 & 0 & 0 & 0 & 0 & 0 & 0 & 0 & 0 & 0\\
\up[v][3] & 0 & 0 & 0 & 0 & 0 & 0 & 0 & 0 & 0 & 0 & 0
\end{array}
\]
(The whole tree has depth $3$, so the $4$th ancestor of every vertex is already above the
root: rows $j{=}2$ and $j{=}3$ are all sentinel.) Now query $\operatorname{lca}(8,11)$. Both
have depth $3$, so $\Delta=0$ and the first loop does nothing; they are unequal. Second loop,
$j=3{:}\ \up[8][3]=\up[11][3]=0$ equal, skip. $j=2{:}\ \up[8][2]=\up[11][2]=0$ equal, skip.
$j=1{:}\ \up[8][1]=2\neq 3=\up[11][1]$, jump:
$u\gets 2,\ v\gets 3$. $j=0{:}\ \up[2][0]=\up[3][0]=1$ equal, skip. Answer:
$\up[u][0]=\up[2][0]=1$. Correct.

Try $\operatorname{lca}(10,7)$ yourself: lift $10$ (depth $3$) up by $\Delta=1$ to its parent
$6$ (depth $2$), then lockstep with $7$ to reach $3$.
\end{labbox}

\begin{trapbox}
\textbf{Three classic binary-lifting bugs.} (1) Forgetting the sentinel: if $\pa(r)$ is left
undefined rather than set to a node $0$ with $\up[0][j]=0$, the recurrence dereferences
garbage when a jump exceeds the root. (2) Equalizing depths using the \emph{wrong} vertex --
you must lift the \emph{deeper} one; swap first. (3) In the second loop, jumping when
$\up[u][j]=\up[v][j]$. That moves both to a common ancestor at or above the LCA and you lose
the answer; the test must be ``jump only when they \emph{differ}.''
\end{trapbox}

\begin{applicationbox}
Binary lifting is not only for LCA. The same $\up$ table answers ``the $k$-th ancestor of
$v$'' in $\Oh(\log n)$ by reading off the bits of $k$, which powers level-ancestor queries,
jump pointers in tries, and the ``walk up to the highest ancestor with property $P$''
searches inside heavy-light decomposition (Chapter 19) and on functional graphs (where each
node has exactly one outgoing edge and $\up[v][j]$ is the $2^j$-th iterate of the successor
function).
\end{applicationbox}

\section{Range Minimum Queries and Sparse Tables}
\label{ch23:rmq}

Before the second LCA method we need its engine. The \emph{range minimum query} problem:
preprocess an array $A[1\ldots m]$ so that for any query interval $[\ell,r]$ we can return
$\min_{\ell\le i\le r} A[i]$ (and, for LCA, the \emph{position} achieving it) quickly.

\begin{intuitionbox}
Minimum has a property that sum does not: it is \emph{idempotent}, $\min(x,x)=x$. So
overlapping ranges may be combined freely -- covering $[\ell,r]$ with two blocks that
\emph{overlap} in the middle changes nothing, because counting an element twice in a $\min$
is harmless. That single fact is what lets RMQ be answered in $\Oh(1)$ after $\Oh(m\log m)$
preprocessing, a luxury sum cannot afford.
\end{intuitionbox}

The \emph{sparse table} precomputes the minimum of every block whose length is a power of
two. Let $\spt[i][j]=\min\{A[i],\dots,A[i+2^j-1]\}$, the minimum of the length-$2^j$ block
starting at $i$. The doubling recurrence mirrors binary lifting: a block of length $2^j$ is
two abutting blocks of length $2^{j-1}$,
\begin{equation}
\spt[i][j] \;=\; \min\bigl(\,\spt[i][j-1],\ \spt[i+2^{j-1}][j-1]\,\bigr).
\label{ch23:spt-recur}
\end{equation}

\begin{algorithm}[H]
\DontPrintSemicolon
\caption{\textsc{BuildSparse}$(A,m)$}
$L \gets \lfloor\log_2 m\rfloor$\;
\lFor{$i \gets 1$ \KwTo $m$}{$\spt[i][0] \gets A[i]$}
\For{$j \gets 1$ \KwTo $L$}{
  \For{$i \gets 1$ \KwTo $m-2^j+1$}{
    $\spt[i][j] \gets \min\bigl(\spt[i][j-1],\ \spt[i+2^{j-1}][j-1]\bigr)$\;
  }
}
\Return $\spt$\;
\end{algorithm}

To answer $[\ell,r]$, let $\textit{len}=r-\ell+1$ and $j=\lfloor\log_2 \textit{len}\rfloor$.
Two blocks of length $2^j$ -- one anchored at $\ell$, one ending at $r$ -- together cover
$[\ell,r]$, since $2^j > \textit{len}/2$ means each block covers more than half the interval
and they overlap in the middle. By idempotency the overlap is harmless:
\begin{equation}
\textsc{Query}(\ell,r) \;=\; \min\bigl(\,\spt[\ell][j],\ \spt[r-2^j+1][j]\,\bigr).
\label{ch23:spt-query}
\end{equation}

\begin{figure}[H]
\centering
\begin{tikzpicture}[>=Stealth,scale=0.92]
\foreach \i/\val in {1/3,2/1,3/4,4/1,5/5,6/9,7/2,8/6,9/5,10/3}{
  \node[draw,minimum width=8mm,minimum height=7mm] (a\i) at (\i,0) {\val};
  \node[font=\scriptsize,below=0.5mm of a\i] {\i};
}
\draw[sourceblue,very thick,decorate,decoration={brace,amplitude=5pt,mirror}]
  ($(a3.south west)+(0,-0.55)$) -- ($(a6.south east)+(0,-0.55)$)
  node[midway,below=4pt,font=\scriptsize]{$\spt[3][2]=\min(4,1,5,9)=1$};
\draw[targetgreen,very thick]
  ($(a5.north west)+(0,0.15)$) -- ($(a8.north east)+(0,0.15)$)
  node[midway,above=2pt,font=\scriptsize]{$\spt[5][2]=\min(5,9,2,6)=2$};
\draw[trapred,thick,->] (3,1.4) -- (8,1.4)
  node[midway,above,font=\scriptsize]{query $[3,8]$, $\textit{len}=6$, $j=2$};
\node[font=\scriptsize,trapred] at (5.5,-1.65)
  {$\min(\spt[3][2],\spt[5][2])=\min(1,2)=1$ \ (blocks overlap on $5,6$)};
\end{tikzpicture}
\caption{A sparse-table query on $A=[3,1,4,1,5,9,2,6,5,3]$. The two length-$4$ blocks
covering $[3,8]$ overlap on positions $5,6$; idempotency of $\min$ makes the double-count
harmless.}
\label{ch23:fig-sparse}
\end{figure}

\begin{whybox}
\textbf{Why two blocks always suffice.} With $j=\lfloor\log_2\textit{len}\rfloor$ we have
$2^j\le\textit{len}<2^{j+1}$, so $2^j>\textit{len}/2$. The first block covers
$[\ell,\ell+2^j-1]$ and the second covers $[r-2^j+1,r]$. Their union starts at $\ell$ and
ends at $r$; they overlap because $\ell+2^j-1 \ge r-2^j+1$ rearranges to
$2\cdot 2^j \ge \textit{len}+1$, which holds since $2^j>\textit{len}/2$ and the two sides
differ by an integer. No gap, hence full coverage, hence the answer is exactly
\eqref{ch23:spt-query}.
\end{whybox}

\begin{complexitybox}
Sparse-table RMQ: $\Th(m\log m)$ preprocessing time and space; $\Oh(1)$ per query (one
lookup of $\lfloor\log_2 \textit{len}\rfloor$, precomputable into a small array, plus one
$\min$). The catch: it is \emph{static}. Updating an array entry can invalidate
$\Th(\log m)$ table cells with no cheap repair; sparse tables are for immutable arrays. For
updates, use a segment tree (Section~\ref{ch23:segtree}). Idempotency is essential -- a
sparse table computes $\min,\max,\gcd$, bitwise-and/or in $\Oh(1)$, but \emph{not} sum,
because double-counting the overlap would corrupt a sum.
\end{complexitybox}

\begin{labbox}
On $A=[3,1,4,1,5,9,2,6,5,3]$ ($m=10$, $L=3$), the $j=1$ row of pairwise minima is
$\spt[\cdot][1]=[1,1,1,1,5,2,2,5,3]$ and the $j=2$ row of length-$4$ minima is
$\spt[\cdot][2]=[1,1,1,1,2,2,2]$. Query $[2,9]$: $\textit{len}=8$, $j=3$,
$\min(\spt[2][3],\spt[2][3])=\spt[2][3]=\min(A[2..9])=1$. Query $[7,10]$: $\textit{len}=4$,
$j=2$, $\min(\spt[7][2],\spt[7][2])=2$. Recompute $[3,8]$ and confirm it equals $1$ as in
Figure~\ref{ch23:fig-sparse}.
\end{labbox}

\section{Euler Tour Plus RMQ: LCA in \texorpdfstring{$\Oh(1)$}{O(1)}}
\label{ch23:euler-lca}

Binary lifting answers in $\Oh(\log n)$. The Euler-tour method answers in $\Oh(1)$ by
reducing LCA to RMQ and dropping a sparse table on top.

\begin{intuitionbox}
Walk the tree with a depth-first traversal, but write down a vertex \emph{every time you
touch it} -- when you first enter, and again each time you return after finishing a child.
This \emph{Euler tour} is a sequence that goes down and up the tree like a heartbeat. Crucial
observation: on the segment of the tour between an occurrence of $u$ and an occurrence of
$v$, you cannot have left the subtree of $\operatorname{lca}(u,v)$, and you must have touched
$\operatorname{lca}(u,v)$ itself (the only way from $u$'s region to $v$'s region is through
their common ancestor). And the LCA is the \emph{shallowest} vertex on that segment --
everything else between them is deeper. So the LCA is the minimum-depth entry in a range of
the tour. That is an RMQ. Figure~\ref{ch23:fig-euler} draws the tour as a depth profile so
this ``shallowest point in the range'' reading is visible at a glance.
\end{intuitionbox}

Concretely, run a DFS from $r$. Maintain three arrays:
\begin{itemize}
\item $\eu[1\ldots 2n-1]$: the vertex visited at each tour step (the \emph{Euler array});
\item $\de[1\ldots 2n-1]$: $\de[k]=\dep(\eu[k])$, the depth at that step;
\item $\fst[v]$: the index of the \emph{first} occurrence of $v$ in $\eu$.
\end{itemize}
A DFS on $T$ produces a tour of exactly $2n-1$ entries: each of the $n-1$ edges is traversed
twice (down and up), giving $2(n-1)$ steps after the root, plus the initial visit to the
root, total $2n-1$.

\begin{algorithm}[H]
\DontPrintSemicolon
\caption{\textsc{EulerTour}$(v)$ -- emit into $\eu,\de$ and record $\fst$}
append $v$ to $\eu$; append $\dep(v)$ to $\de$; let $k$ be the new last index\;
\lIf{$\fst[v]$ undefined}{$\fst[v] \gets k$}
\For{each child $c$ of $v$}{
  \textsc{EulerTour}$(c)$\;
  append $v$ to $\eu$; append $\dep(v)$ to $\de$ \tcp*{re-touch $v$ on the way back up}
}
\end{algorithm}

\begin{theorem}[LCA is RMQ on the Euler tour]
\label{ch23:thm-euler}
For any $u,v$, let $a=\fst[u]$ and $b=\fst[v]$ with (without loss of generality) $a\le b$.
Let $k^\star$ be an index in $[a,b]$ minimizing $\de[\cdot]$. Then
$\eu[k^\star]=\operatorname{lca}(u,v)$.
\end{theorem}

\begin{proof}
Let $\ell=\operatorname{lca}(u,v)$. Two claims.

\emph{(i) An occurrence of $\ell$ lies in $[a,b]$.} The tour enters the subtree of $\ell$
exactly once (at $\ell$'s first occurrence) and leaves it once, and between those moments it
visits every descendant of $\ell$ and only those, in one contiguous block of tour-time. Both
$u$ and $v$ are descendants of $\ell$, so their first occurrences $a,b$ both lie inside that
block, which begins and ends at occurrences of $\ell$. To travel in the tour from the visit
of $u$ to the visit of $v$, the walk must climb to $\ell$ and descend again, because $u$ and
$v$ sit in different child-subtrees of $\ell$ (or one is $\ell$); either way some occurrence
of $\ell$ falls in $[a,b]$. Hence $\ell=\eu[k]$ for some $k\in[a,b]$, with
$\de[k]=\dep(\ell)$.

\emph{(ii) No vertex in $[a,b]$ is shallower than $\ell$.} Every vertex touched between
times $a$ and $b$ is a descendant of $\ell$ (the walk never leaves $\ell$'s subtree in that
interval, by the contiguity in (i)). Descendants of $\ell$ have depth $\ge\dep(\ell)$, with
equality only for $\ell$ itself. So $\min_{k\in[a,b]}\de[k]=\dep(\ell)$, attained exactly at
occurrences of $\ell$. Therefore any minimizing index $k^\star$ has
$\eu[k^\star]=\ell$.\qedhere
\end{proof}

The recipe: build $\eu,\de,\fst$ in one DFS ($\Oh(n)$), build a sparse table over $\de$
($\Oh(n\log n)$), and answer each $\operatorname{lca}(u,v)$ by an RMQ on
$[\min(\fst[u],\fst[v]),\,\max(\fst[u],\fst[v])]$ in $\Oh(1)$, returning the \emph{vertex}
$\eu[\cdot]$ at the minimizing index.

\begin{labbox}
\textbf{Euler tour of $T$.} A DFS visiting children in increasing order produces (with root
$1$):
\[
\renewcommand{\arraystretch}{1.15}
\begin{array}{c|ccccccccccccccccccccc}
k    & 1&2&3&4&5&6&7&8&9&10&11&12&13&14&15&16&17&18&19&20&21\\\hline
\eu  & 1&2&4&2&5&8&5&9&5&2 &1 &3 &6 &10&6 &11&6 &3 &7 &3 &1\\
\de  & 0&1&2&1&2&3&2&3&2&1 &0 &1 &2 &3 &2 &3 &2 &1 &2 &1 &0
\end{array}
\]
($2n-1=21$ entries, good.) First occurrences, for $v=1\ldots 11$:
$\fst=[1,2,12,3,5,13,19,6,8,14,16]$.

Query $\operatorname{lca}(8,11)$: $\fst[8]=6$, $\fst[11]=16$, so RMQ on
$\de[6\ldots16]=[3,2,3,2,1,0,1,2,3,2,3]$. The minimum is $0$ at $k=11$, and $\eu[11]=1$.
Answer $1$ -- matches binary lifting.

Query $\operatorname{lca}(10,7)$: $\fst[10]=14$, $\fst[7]=19$, RMQ on
$\de[14\ldots19]=[3,2,3,2,1,2]$, minimum $1$ at $k=18$, $\eu[18]=3$. Answer $3$. Correct.
\end{labbox}

\begin{figure}[H]
\centering
\begin{tikzpicture}[>=Stealth,xscale=0.6,yscale=0.9]
\def\dseq{0,1,2,1,2,3,2,3,2,1,0,1,2,3,2,3,2,1,2,1,0}
\def\vseq{1,2,4,2,5,8,5,9,5,2,1,3,6,10,6,11,6,3,7,3,1}
% axis baseline
\draw[edgegray,->] (0.4,0) -- (21.6,0) node[right,font=\scriptsize]{$k$};
% depth profile (heartbeat)
\foreach \d [count=\k from 1] in \dseq {
  \node[circle,fill=nodeblue,inner sep=1.1pt] (p\k) at (\k,-\d) {};
}
\foreach \k in {1,...,20}{\pgfmathtruncatemacro{\kk}{\k+1}\draw[nodeblue,thick] (p\k) -- (p\kk);}
% label the vertices along the curve
\foreach \v [count=\k from 1] in \vseq {
  \node[font=\scriptsize,below=0.6pt of p\k]{\v};
}
% highlight the query range [6,16] for lca(8,11)
\draw[trapred,very thick,decorate,decoration={brace,amplitude=4pt}]
  (6,0.55) -- (16,0.55) node[midway,above=3pt,font=\scriptsize]{range $[\fst[8],\fst[11]]=[6,16]$};
\node[circle,draw=trapred,very thick,inner sep=2.4pt] at (11,0){};
\node[font=\scriptsize,trapred,above right=1pt and 1pt of p11]{shallowest $=1$};
\end{tikzpicture}
\caption{The Euler tour of $T$ drawn as a depth profile (down is deeper). Each dot is a tour
step $k$, plotted at height $-\de[k]$ and labelled with $\eu[k]$. The query
$\operatorname{lca}(8,11)$ scans the range $[6,16]$; its shallowest (topmost) point sits at
$k=11$, where $\eu[11]=1$ -- the LCA.}
\label{ch23:fig-euler}
\end{figure}

\begin{trapbox}
\textbf{Query the right range.} You must use \emph{first} occurrences and order them with
$\min/\max$; mixing a first occurrence of one vertex with a \emph{later} occurrence of the
other can skip past the LCA. Using $\fst[u]$ and $\fst[v]$ (sorted, in either original order)
is what the proof relies on. Note too that the Euler array here has length $2n-1$; a common
variant emits a vertex on \emph{every} edge traversal with a trailing copy, or uses
entry/exit timestamps, landing at $2n$ entries -- pick one convention and never mix them.
\end{trapbox}

\begin{historybox}
The reduction of LCA to RMQ, and an $\langle\Oh(n),\Oh(1)\rangle$ algorithm for it, is due to
Harel and Tarjan (1984), who first achieved constant-time LCA after linear preprocessing.
Their method was intricate; Bender and Farach-Colton (2000) gave the clean modern version:
reduce LCA to RMQ via the Euler tour, then exploit that consecutive depths differ by exactly
$\pm 1$ -- the \emph{$\pm1$ RMQ} -- to reach true $\Oh(n)$ preprocessing by blocking the array
and tabulating the $\Oh(\sqrt{m})$ distinct small blocks. The sparse-table version we built
gives $\langle\Oh(n\log n),\Oh(1)\rangle$, which is what nearly everyone actually codes.
\end{historybox}

\begin{aestheticbox}
The $\pm1$ structure is the secret to linear preprocessing. Because adjacent depths in $\de$
differ by exactly $1$, a block of size $b=\tfrac12\log_2 m$ is determined by its first value
plus a $\pm1$ pattern of $b-1$ bits, so there are only $2^{b-1}=\Oh(\sqrt m)$ distinct block
shapes. Precompute the in-block RMQ for all of them once, store one block-minimum per block
in a sparse table over $\Oh(m/\log m)$ blocks, and the total preprocessing collapses to
$\Oh(m)$ with $\Oh(1)$ queries. This ``few distinct small pieces, tabulate them all'' move is
the Four-Russians idea, and it recurs across algorithm design.
\end{aestheticbox}

\section{Segment Trees: Queries with Updates}
\label{ch23:segtree}

The sparse table is frozen. The moment the array can change, we want a \emph{segment tree}:
$\Oh(\log m)$ for both point updates and range queries, and -- with lazy propagation -- range
updates too. It is the universal range structure, working for any associative combine
operation (sum, min, max, gcd, matrix product), idempotent or not.

\begin{intuitionbox}
Recursively cut the array in half. The root owns $[1,m]$; its children own the two halves;
leaves own single cells. Each internal node stores the combined value (say sum, or min) of
its segment. A query interval $[\ell,r]$ shatters into $\Oh(\log m)$ \emph{canonical}
segments -- whole nodes that tile $[\ell,r]$ -- and we combine their stored values. An update
to one cell touches only the $\Oh(\log m)$ nodes on the root-to-leaf path above it.
\end{intuitionbox}

\begin{figure}[H]
\centering
\begin{tikzpicture}[every node/.style={draw,rounded corners,font=\scriptsize,
  minimum height=5mm,inner sep=2pt},level distance=11mm,>=Stealth,
  level 1/.style={sibling distance=46mm},
  level 2/.style={sibling distance=23mm},
  level 3/.style={sibling distance=12mm}]
\node[fill=blue!6] {$[1,8]{:}31$}
 child {node[fill=blue!6] {$[1,4]{:}9$}
   child {node {$[1,2]{:}4$}
     child {node[fill=green!10] {$[1]{:}3$}}
     child {node[fill=green!10] {$[2]{:}1$}}}
   child {node[fill=orange!12] {$[3,4]{:}5$}
     child {node {$[3]{:}4$}}
     child {node {$[4]{:}1$}}}}
 child {node[fill=blue!6] {$[5,8]{:}22$}
   child {node[fill=orange!12] {$[5,6]{:}14$}
     child {node {$[5]{:}5$}}
     child {node {$[6]{:}9$}}}
   child {node {$[7,8]{:}8$}
     child {node {$[7]{:}2$}}
     child {node {$[8]{:}6$}}}};
\node[draw=none,font=\scriptsize\itshape,text width=33mm] at (6.4,-1.1)
  {\textbf{sum} query $[3,6]$ splits into the two shaded canonical nodes
  $[3,4]{=}5$ and $[5,6]{=}14$; answer $19$.};
\end{tikzpicture}
\caption{Segment tree over $A=[3,1,4,1,5,9,2,6]$ storing subtree sums. The interval $[3,6]$
decomposes into two canonical nodes (orange).}
\label{ch23:fig-segtree}
\end{figure}

We store the tree in a flat array $t[1\ldots 4m]$: node $1$ is the root, node $p$'s children
are $2p$ and $2p{+}1$. (Size $4m$ safely covers any $m$; the tree has at most $2\cdot
2^{\lceil\log_2 m\rceil}\le 4m$ nodes.) Let $\oplus$ be the combine operator with identity
$e$ (e.g. $0$ for sum, $+\infty$ for min). Build recursively:

\begin{algorithm}[H]
\DontPrintSemicolon
\caption{\textsc{Build}$(p,\ell,r)$ -- node $p$ owns $[\ell,r]$}
\If{$\ell = r$}{$t[p] \gets A[\ell]$\; \Return}
$\textit{mid} \gets \lfloor(\ell+r)/2\rfloor$\;
\textsc{Build}$(2p,\ell,\textit{mid})$\; \textsc{Build}$(2p{+}1,\textit{mid}{+}1,r)$\;
$t[p] \gets t[2p] \oplus t[2p{+}1]$\;
\end{algorithm}

\begin{algorithm}[H]
\DontPrintSemicolon
\caption{\textsc{Query}$(p,\ell,r,\, q_\ell,q_r)$ -- combine $A[q_\ell\ldots q_r]$}
\lIf{$q_r < \ell$ \textbf{or} $r < q_\ell$}{\Return $e$ \tcp*{node disjoint from query}}
\lIf{$q_\ell \le \ell$ \textbf{and} $r \le q_r$}{\Return $t[p]$ \tcp*{node inside query}}
$\textit{mid} \gets \lfloor(\ell+r)/2\rfloor$\;
\Return $\textsc{Query}(2p,\ell,\textit{mid},q_\ell,q_r) \,\oplus\, \textsc{Query}(2p{+}1,\textit{mid}{+}1,r,q_\ell,q_r)$\;
\end{algorithm}

\begin{algorithm}[H]
\DontPrintSemicolon
\caption{\textsc{Update}$(p,\ell,r,\, i,\, x)$ -- set $A[i]\gets x$}
\If{$\ell = r$}{$t[p] \gets x$\; \Return}
$\textit{mid} \gets \lfloor(\ell+r)/2\rfloor$\;
\lIf{$i \le \textit{mid}$}{\textsc{Update}$(2p,\ell,\textit{mid},i,x)$}
\lElse{\textsc{Update}$(2p{+}1,\textit{mid}{+}1,r,i,x)$}
$t[p] \gets t[2p] \oplus t[2p{+}1]$\;
\end{algorithm}

\begin{whybox}
\textbf{Why a query touches only $\Oh(\log m)$ nodes.} Call a node \emph{partial} if the
query interval cuts through its segment without containing it (the recursion continues),
\emph{total} if its segment is contained (returns $t[p]$, no recursion), \emph{null} if
disjoint (returns $e$). Recursion continues only at partial nodes. Claim: at each depth of
the tree, at most two nodes are partial. The query interval $[q_\ell,q_r]$ has a left endpoint
and a right endpoint; a node at a given depth is partial only if one of these endpoints lies
strictly inside its segment. The segments at one depth are disjoint, so $q_\ell$ falls in at
most one and $q_r$ in at most one -- at most two partial nodes per depth. Each partial node
does $\Oh(1)$ work and examines at most two children, so the total node count is
$\Oh(\text{depth})=\Oh(\log m)$. \textsc{Update} follows a single root-to-leaf path, also
$\Oh(\log m)$.
\end{whybox}

\subsection{Lazy propagation: range updates}

A point update is easy; a \emph{range} update -- ``add $x$ to every $A[i]$ with
$\ell\le i\le r$'' -- would naively touch $\Th(r-\ell)$ leaves. \emph{Lazy propagation} fixes
this: when an update fully covers a node's segment, apply it to the node's aggregate and park
a \emph{pending tag} on the node instead of recursing. The tag is pushed down to children
only when a later operation needs to descend through the node.

\begin{algorithm}[H]
\DontPrintSemicolon
\caption{\textsc{RangeAdd}$(p,\ell,r,\, q_\ell,q_r,\, x)$ -- add $x$ to $A[q_\ell\ldots q_r]$}
\lIf{$q_r < \ell$ \textbf{or} $r < q_\ell$}{\Return}
\If{$q_\ell \le \ell$ \textbf{and} $r \le q_r$}{
  $t[p] \mathrel{+}= (r-\ell+1)\cdot x$\; $\lz[p] \mathrel{+}= x$\; \Return
}
\textsc{PushDown}$(p,\ell,r)$\;
$\textit{mid} \gets \lfloor(\ell+r)/2\rfloor$\;
\textsc{RangeAdd}$(2p,\ell,\textit{mid},q_\ell,q_r,x)$\;
\textsc{RangeAdd}$(2p{+}1,\textit{mid}{+}1,r,q_\ell,q_r,x)$\;
$t[p] \gets t[2p] + t[2p{+}1]$\;
\end{algorithm}

\textsc{PushDown}$(p,\ell,r)$ takes the tag $\lz[p]$, adds $\lz[p]\cdot(\text{child size})$
to each child's $t$ and $\lz[p]$ to each child's tag, then clears $\lz[p]$. Because a tag on a
fully covered node already reflects the correct aggregate in $t[p]$, and is pushed only when
we descend, the invariant ``$t[p]$ is correct \emph{including} the effect of $\lz[p]$, while
$\lz[p]$ holds updates not yet reflected in the children'' is maintained throughout. Range
update and range query are each $\Oh(\log m)$ by the same two-partial-nodes-per-depth
argument: pushing down adds $\Oh(1)$ per visited node.

\begin{trapbox}
\textbf{The lazy invariant is delicate.} Two ways to corrupt it: (1) querying or updating a
child without \textsc{PushDown}-ing the parent first -- the child still hides a stale tag;
(2) forgetting that a node's stored aggregate $t[p]$ must \emph{already include} its own tag.
Rule of thumb: a node's $t[p]$ is always correct; a node's $\lz[p]$ is debt owed to its
children, paid only on the way down. For non-commuting tag compositions (range \emph{assign}
mixed with range \emph{add}) the order in which composed tags apply matters; get it wrong and
the bug is maddening to find.
\end{trapbox}

\begin{complexitybox}
Segment tree: $\Th(m)$ build, $\Oh(\log m)$ per point update / range query; with lazy
propagation $\Oh(\log m)$ per range update too. Space $\Th(m)$. Works for any associative
$\oplus$. As an RMQ structure it gives $\langle\Oh(m),\Oh(\log m)\rangle$ -- slower queries
than a sparse table, but it supports updates, which the sparse table cannot.
\end{complexitybox}

\section{Fenwick (Binary Indexed) Trees}
\label{ch23:fenwick}

For the common special case of \emph{prefix-sum} queries with point updates, the Fenwick
tree is a featherweight alternative: a single array, a handful of bit-twiddling lines,
$\Oh(\log m)$ per operation, and a constant factor that beats the segment tree.

\begin{intuitionbox}
Index $i$ in binary ends in some run of zeros; let $\lowbit(i)=i\,\&\,(-i)$ be the value of
its lowest set bit. Fenwick cell $F[i]$ stores the sum of the $\lowbit(i)$ array elements
ending at $i$, i.e. $A[i-\lowbit(i)+1\ldots i]$. Then any prefix sum $A[1\ldots i]$ is
assembled by peeling off lowest bits: $F[i]$ covers a suffix ending at $i$, then jump to
$i-\lowbit(i)$ and add its block, and so on down to $0$. The binary expansion of $i$ has at
most $\log_2 m$ ones, so a prefix sum is at most $\log_2 m$ additions.
\end{intuitionbox}

\begin{figure}[H]
\centering
\begin{tikzpicture}[>=Stealth,scale=0.95]
\foreach \i in {1,...,8}{
  \node[draw,minimum size=7mm] (f\i) at (\i,0) {$F_{\i}$};
  \node[font=\scriptsize,below=0.5mm of f\i] {\i};
}
\node[font=\scriptsize,sourceblue] at (4.5,1.05)
  {$F[i]$ covers the block $A[i-\lowbit(i)+1\,..\,i]$};
\node[font=\scriptsize] at (4.5,-1.0)
  {prefix$(7)=F[7]+F[6]+F[4]$\quad(indices $7\to6\to4\to0$, peeling lowest bits)};
\end{tikzpicture}
\caption{Fenwick coverage. $F[8]$ covers all of $A[1..8]$, $F[4]$ covers $A[1..4]$,
$F[6]$ covers $A[5..6]$, $F[7]$ covers just $A[7]$.}
\label{ch23:fig-fenwick}
\end{figure}

\begin{algorithm}[H]
\DontPrintSemicolon
\caption{Fenwick operations on $F[1\ldots m]$ (all zero initially)}
\SetKwProg{Fn}{function}{:}{}
\Fn{\textsc{Add}$(i,x)$ \tcp*{$A[i] \mathrel{+}= x$}}{
  \While{$i \le m$}{$F[i] \mathrel{+}= x$\; $i \gets i + \lowbit(i)$\;}
}
\Fn{\textsc{PrefixSum}$(i)$ \tcp*{returns $A[1]+\dots+A[i]$}}{
  $s \gets 0$\;
  \While{$i > 0$}{$s \mathrel{+}= F[i]$\; $i \gets i - \lowbit(i)$\;}
  \Return $s$\;
}
\end{algorithm}

A range sum $A[\ell\ldots r]$ is $\textsc{PrefixSum}(r)-\textsc{PrefixSum}(\ell-1)$. The two
loops walk \emph{opposite} bit-jumps: queries strip the lowest set bit ($i \to i-\lowbit(i)$),
updates climb to the next cell that contains $i$ ($i \to i+\lowbit(i)$). Each visits at most
$\Oh(\log m)$ cells because each step clears or sets a distinct bit.

\begin{whybox}
\textbf{Why \textsc{Add} hits exactly the cells containing $i$.} Cell $F[j]$ contains $A[i]$
iff $j-\lowbit(j) < i \le j$. The smallest such $j\ge i$ is $i$ itself; the next is obtained
by $j \mapsto j+\lowbit(j)$, which jumps to the smallest larger index whose block reaches
back to $i$, and iterating gives precisely the increasing chain of containing cells, ending
once it exceeds $m$. Dually, $i \mapsto i-\lowbit(i)$ in \textsc{PrefixSum} partitions $[1,i]$
into the disjoint blocks $F[i],F[i-\lowbit(i)],\dots$ whose lengths are the powers of two in
the binary representation of $i$, so their union is exactly $[1,i]$ and their sum is the
prefix sum.
\end{whybox}

\begin{labbox}
\textbf{Trace on $A=[3,1,4,1,5,9,2,6]$.} Building by \textsc{Add}$(i,A[i])$ for each $i$
yields $F=[3,4,4,9,5,14,2,31]$. Check: $F[8]=A[1..8]=31$? The sum is
$3{+}1{+}4{+}1{+}5{+}9{+}2{+}6=31$, yes. $F[6]=A[5]{+}A[6]=14$, $F[7]=A[7]=2$,
$F[4]=A[1..4]=9$. \textsc{PrefixSum}$(7)=F[7]+F[6]+F[4]=2+14+9=25$, and indeed $A[1..7]=25$.
Now \textsc{Add}$(3,+10)$ visits $3\to4\to8$: $F[3]{=}14,F[4]{=}19,F[8]{=}41$. Re-query
$[3,6]=\textsc{PrefixSum}(6)-\textsc{PrefixSum}(2)=(F[6]{+}F[4])-(F[2])=(14{+}19)-4=29$,
matching $A[3..6]$ after the update: $14{+}1{+}5{+}9=29$. Correct.
\end{labbox}

\begin{trapbox}
\textbf{Fenwick is 1-indexed.} The jumps $i\pm\lowbit(i)$ break at $i=0$ ($\lowbit(0)=0$
loops forever); always index from $1$, and let \textsc{PrefixSum} stop at $i=0$. A second
trap: a plain Fenwick tree handles \emph{prefix-decomposable} aggregates -- sums, xor,
counts. It does \emph{not} handle range \emph{min}, because min has no inverse, so you cannot
subtract one prefix from another. (Prefix-min from index $1$ is fine; arbitrary range-min
needs a sparse table or segment tree.)
\end{trapbox}

\begin{complexitybox}
Fenwick: $\Th(m)$ build (or $\Th(m)$ in-place), $\Oh(\log m)$ point update and prefix/range
sum, $\Th(m)$ space with the smallest constant of any structure here. Range-update /
point-query and range-update / range-query variants exist via one or two Fenwick trees over
difference arrays, still $\Oh(\log m)$.
\end{complexitybox}

\begin{applicationbox}
Fenwick trees are the standard way to count \emph{inversions} (pairs $i<j$ with $A[i]>A[j]$):
sweep left to right, query how many already-seen values exceed the current one. Combined with
the Euler tour of Section~\ref{ch23:euler-lca}, a Fenwick tree over entry/exit timestamps
answers subtree-sum and subtree-update queries -- the bridge between this chapter and
Chapter 21's Euler-tour-on-subtrees technique.
\end{applicationbox}

\section{The LCA \texorpdfstring{$\leftrightarrow$}{<->} RMQ Equivalence}
\label{ch23:equivalence}

We reduced LCA to RMQ (Theorem~\ref{ch23:thm-euler}). The reduction runs the other way too,
and the round trip reveals that the two problems are, up to linear-time transformation, the
\emph{same}. This is more than a curiosity: a result for either problem transfers to the
other, and it explains why both admit $\langle\Oh(n),\Oh(1)\rangle$ solutions.

\begin{theorem}[RMQ reduces to LCA via the Cartesian tree]
\label{ch23:thm-cartesian}
Given an array $A[1\ldots m]$, build its \emph{Cartesian tree} $C$: the root is the position
of the global minimum; its left subtree is the Cartesian tree of the elements to the left of
that minimum, its right subtree that of the elements to the right (ties broken, say,
leftmost). Then for any $\ell\le r$, the position of the minimum of $A[\ell\ldots r]$ equals
$\operatorname{lca}_C(\ell,r)$.
\end{theorem}

\begin{proof}
The Cartesian tree has the heap property on values (a node's value is $\le$ its descendants',
since each subtree is built around its own minimum) and the binary-search-tree property on
\emph{positions} (the left subtree holds smaller indices, the right subtree larger). Take
$\ell,r$ and let $w=\operatorname{lca}_C(\ell,r)$ sit at position $p$. By the BST property
$p\in[\ell,r]$: otherwise both $\ell$ and $r$ would lie in the same child subtree of $w$,
contradicting $w$ being their LCA. By the heap property, $A[p]$ is $\le$ every value in the
subtree rooted at $w$, which contains all positions in $[\ell,r]$ (they all descend from
their LCA). Hence $A[p]=\min A[\ell\ldots r]$, and $p$ is a minimizing position.
\end{proof}

\begin{figure}[H]
\centering
\begin{tikzpicture}[every node/.style={circle,draw=structurepurple,thick,fill=purple!4,
  minimum size=8mm,inner sep=0pt,font=\scriptsize},level distance=11mm,>=Stealth,
  level 1/.style={sibling distance=26mm},
  level 2/.style={sibling distance=14mm},
  level 3/.style={sibling distance=10mm}]
\node {$A_2{:}1$}
 child {node {$A_1{:}3$}}
 child {node {$A_4{:}1$}
   child {node {$A_3{:}4$}}
   child {node {$A_7{:}2$}
     child {node {$A_5{:}5$}
       child {node[draw=none,fill=none]{}}
       child {node {$A_6{:}9$}}}
     child {node {$A_8{:}6$}}}};
\node[draw=none,fill=none,font=\scriptsize\itshape,text width=34mm] at (5.3,-1.5)
  {Cartesian tree of $A=[3,1,4,1,5,9,2,6]$. $\operatorname{lca}(5,8)$ is the node
  $A_7{:}2$, and $\min A[5..8]=2$ at position $7$.};
\end{tikzpicture}
\caption{Cartesian tree: RMQ on $A$ becomes LCA on $C$. Node labels are
position${:}$value; the tree is min-heap-ordered on value and BST-ordered on position.}
\label{ch23:fig-cartesian}
\end{figure}

The Cartesian tree of $A[1\ldots m]$ is built in $\Oh(m)$ by a single left-to-right scan with
a stack: maintain the right spine of the tree built so far; to insert $A[i]$, pop spine
elements whose value exceeds $A[i]$ (they become $A[i]$'s left subtree), then attach $A[i]$ as
the right child of whatever remains on top.

\begin{algorithm}[H]
\DontPrintSemicolon
\caption{\textsc{BuildCartesian}$(A,m)$ -- $\Oh(m)$ stack scan}
$\textit{stack} \gets$ empty\;
\For{$i \gets 1$ \KwTo $m$}{
  $\textit{last} \gets \textsc{nil}$\;
  \While{$\textit{stack}$ nonempty \textbf{and} $A[\textit{stack.top}] > A[i]$}{
    $\textit{last} \gets \textsc{pop}(\textit{stack})$\;
  }
  \lIf{$\textit{last} \neq \textsc{nil}$}{$\textsc{leftchild}(i) \gets \textit{last}$}
  \lIf{$\textit{stack}$ nonempty}{$\textsc{rightchild}(\textit{stack.top}) \gets i$}
  \textsc{push}(\textit{stack}, $i$)\;
}
\Return root $=$ bottom of $\textit{stack}$\;
\end{algorithm}

Each index is pushed and popped at most once, so the build is $\Oh(m)$ amortized.

\begin{whybox}
\textbf{Closing the loop.} Put the two reductions together. \emph{LCA $\to$ RMQ}: the Euler
tour turns a tree on $n$ nodes into a $\pm1$ array of length $2n-1$ ($\Oh(n)$).
\emph{RMQ $\to$ LCA}: the Cartesian tree turns an array of length $m$ into a tree on $m$
nodes ($\Oh(m)$). So an $\langle f(n),g(n)\rangle$ algorithm for either problem yields an
$\langle \Oh(n)+f(\Oh(n)),\,\Oh(1)+g(\Oh(n))\rangle$ algorithm for the other. Plug in the
$\pm1$-RMQ $\langle\Oh(n),\Oh(1)\rangle$ result to get $\langle\Oh(n),\Oh(1)\rangle$
\emph{LCA} -- and then, via Cartesian trees, $\langle\Oh(n),\Oh(1)\rangle$ \emph{general}
RMQ. The two problems stand or fall together.
\end{whybox}

\begin{connectionbox}
This equivalence is a recurring shape in this book. The Euler tour reappears in Chapter 21 as
the engine for subtree and path queries; binary lifting is the static cousin of the link-cut
trees of Chapter 22; and the heavy-light decomposition of Chapter 19 uses exactly the segment
trees built here to answer path-aggregate queries. On the spectral side, the Laplacian
$\Lap=\mathbf{D}-\Adjm$ of a tree has quadratic form
$x^\top \Lap x = \sum_{(u,v)\in E}(x_u-x_v)^2$, which telescopes along root-to-LCA paths, so
LCA is the combinatorial skeleton under electrical-flow computations on trees -- a bridge to
the Laplacian-solver material later in this book and in the companion linear-algebra volume.
\end{connectionbox}

\section{Choosing a Structure}
\label{ch23:choosing}

\begin{center}
\renewcommand{\arraystretch}{1.3}
\begin{tabular}{@{}lllll@{}}
\toprule
Structure & Preprocess & Query & Update & Operations\\
\midrule
Sparse table        & $\Oh(m\log m)$ & $\Oh(1)$       & --                  & idempotent (min/max/gcd)\\
Segment tree        & $\Oh(m)$       & $\Oh(\log m)$  & $\Oh(\log m)$       & any associative $\oplus$\\
Segment tree + lazy & $\Oh(m)$       & $\Oh(\log m)$  & $\Oh(\log m)$ range & associative + tags\\
Fenwick tree        & $\Oh(m)$       & $\Oh(\log m)$  & $\Oh(\log m)$       & invertible (sum/xor)\\
Binary-lifting LCA  & $\Oh(n\log n)$ & $\Oh(\log n)$  & --                  & LCA, $k$-th ancestor\\
Euler + sparse LCA  & $\Oh(n\log n)$ & $\Oh(1)$       & --                  & LCA (static tree)\\
\bottomrule
\end{tabular}
\end{center}

\begin{focusbox}
The decision tree in practice: \emph{Static array, idempotent query, want $\Oh(1)$?} sparse
table. \emph{Updates needed, or non-idempotent combine?} segment tree (add lazy tags for
range updates). \emph{Only prefix/range sums with point updates?} Fenwick -- least code,
smallest constant. \emph{LCA with an $\Oh(\log n)$ budget and you also want $k$-th ancestor?}
binary lifting. \emph{LCA queries dominate and the tree is fixed?} Euler tour plus sparse
table for $\Oh(1)$.
\end{focusbox}

\section{Exercises}
\label{ch23:exercises}

\begin{exercise}[Distance and the LCA]
\label{ch23:ex-distance}
In a tree with nonnegative edge weights, let $D(v)$ be the weighted distance from the root to
$v$. Prove that the weighted distance between $u$ and $v$ equals
$D(u)+D(v)-2D(\operatorname{lca}(u,v))$. Then describe an $\Oh(\log n)$-per-query algorithm
combining binary lifting with a precomputed $D$ array.
\end{exercise}

\begin{exercise}[Implement binary lifting]
\label{ch23:ex-impl-lift}
Implement \textsc{BuildLifting} and \textsc{LiftLCA} and test them against \textsc{NaiveLCA}
on random trees of up to $10^5$ nodes. Add a function that returns the $k$-th ancestor of a
node, or the root's sentinel if $k>\dep(v)$, using only the $\up$ table.
\end{exercise}

\begin{exercise}[Break the lockstep]
\label{ch23:ex-break}
A student writes the second loop of \textsc{LiftLCA} to jump when $\up[u][j]=\up[v][j]$ (the
negation of the correct test). Exhibit a small tree and a query where this returns a wrong
answer, and say in one sentence what the buggy loop actually computes.
\end{exercise}

\begin{exercise}[Sparse table for gcd]
\label{ch23:ex-gcd}
Argue that $\gcd$ is associative and idempotent, so a sparse table answers range-gcd in
$\Oh(1)$ comparisons. If entries are big integers of $B$ bits, what is the true per-query
time, and how does that change the bound? Contrast with range-\emph{sum}, and explain
precisely why the overlapping-block trick fails for sums.
\end{exercise}

\begin{exercise}[Euler-tour length and bracketing]
\label{ch23:ex-euler-len}
Prove that the Euler tour produced by \textsc{EulerTour} has exactly $2n-1$ entries for a tree
on $n$ nodes, by a counting argument over edge traversals. Then prove that the first and last
occurrences of any vertex $v$ bracket exactly the tour steps spent inside $v$'s subtree.
\end{exercise}

\begin{exercise}[Lazy assign \emph{and} add]
\label{ch23:ex-lazy}
Extend the lazy segment tree to support two range updates -- ``assign $x$ to $A[\ell\ldots r]$''
and ``add $x$ to $A[\ell\ldots r]$,'' interleaved arbitrarily -- with range-sum queries.
Define the tag type and its composition rule, and explain why the order of applying a pending
assign versus a pending add matters. Give the \textsc{PushDown} that respects it.
\end{exercise}

\begin{exercise}[Inversions with Fenwick]
\label{ch23:ex-inv}
Using a Fenwick tree, count the inversions of a permutation of $1\ldots m$ in $\Oh(m\log m)$.
Then modify your code to compute, for each $i$, the number of $j>i$ with $A[j]<A[i]$ (the
per-element inversion count), still in $\Oh(m\log m)$.
\end{exercise}

\begin{exercise}[Cartesian tree by hand]
\label{ch23:ex-cartesian}
Build the Cartesian tree of $A=[5,2,8,1,9,3,7,4]$ using the stack scan, drawing the right
spine after each insertion. Verify on three ranges that $\operatorname{lca}_C(\ell,r)$ is the
position of $\min A[\ell\ldots r]$.
\end{exercise}

\begin{exercise}[$\pm1$ RMQ in $\Oh(n)$ preprocessing]
\label{ch23:ex-pm1}
Implement the Bender--Farach-Colton $\pm1$ RMQ: split the depth array into blocks of size
$b=\tfrac12\log_2 m$, build a sparse table over the $\Oh(m/b)$ block minima, and tabulate
in-block answers for all $\Oh(\sqrt m)$ distinct $\pm1$ patterns. Prove that the total
preprocessing is $\Oh(n)$ and queries are $\Oh(1)$, and use it to get linear-preprocessing
LCA.
\end{exercise}

\begin{exercise}[Offline LCA, Tarjan]
\label{ch23:ex-tarjan}
Look up Tarjan's \emph{offline} LCA algorithm, which answers a batch of queries during a
single DFS using a union-find structure. Implement it, and prove its running time is
$\Oh((n+q)\,\alpha(n))$ where $\alpha$ is the inverse Ackermann function. When is the offline
method preferable to the online structures of this chapter?
\end{exercise}

\begin{exercise}[Persistent segment tree]
\label{ch23:ex-persistent}
A \emph{persistent} segment tree keeps every past version by copying only the $\Oh(\log m)$
nodes on an update path. Implement one supporting point updates, and use $m$ stacked versions
to answer ``the $k$-th smallest value in $A[\ell\ldots r]$'' (a static range-$k$th query) in
$\Oh(\log m)$ per query. Explain why the total space is $\Oh(m\log m)$.
\end{exercise}
